\phantomsection
\chapter*{Alana}
\addcontentsline{toc}{section}{Alana --- Nisa Nurfadhillah}
\penulis{Nisa Nurfadhillah}


Di sebuah desa kecil yang dikelilingi sawah dan perbukitan, tinggallah seorang gadis bernama Alana bersama ibunya. Kehidupan mereka sangat sederhana. Mereka tinggal di sebuah rumah kecil yang jauh dari kemewahan. Meskipun begitu, Alana tumbuh sebagai anak yang ceria, rajin, dan menyayangi ibunya.

Sejak kecil, Alana tidak pernah benar-benar mengenal sosok ayah. Ia hanya mengetahui bahwa ayahnya masih hidup dan tinggal jauh dari desa. Ibunya selalu menyimpan kisah tentang ayah Alana dengan rapat. Tidak banyak yang diketahui Alana mengenai alasan mereka tidak hidup bersama.

Alana terbiasa menjalani kehidupan sederhana. Setiap pagi ia berangkat sekolah dengan seragam yang telah berkali-kali dicuci. Sepulang sekolah, ia membantu ibunya mengerjakan pekerjaan rumah. Pada malam hari, mereka menghabiskan waktu bersama dalam rumah kecil yang sederhana.

Bagi Alana, kebahagiaan tidak pernah ditentukan oleh banyaknya harta. Kehadiran ibunya sudah lebih dari cukup.

Namun, kebahagiaan itu tidak berlangsung selamanya.

Suatu hari, ibu Alana jatuh sakit. Kondisinya semakin memburuk hingga membuat Alana harus menghadapi kenyataan yang paling ia takuti. Setelah beberapa waktu berjuang melawan penyakit, ibunya akhirnya meninggal dunia.

Kepergian tersebut membuat kehidupan Alana berubah sepenuhnya.

Rumah kecil yang sebelumnya dipenuhi kehangatan mendadak terasa begitu sepi. Tidak ada lagi seseorang yang menunggunya pulang. Tidak ada lagi makanan sederhana yang disiapkan dengan penuh kasih sayang. Tidak ada lagi sosok yang selalu menjadi tempat Alana berbagi cerita.

Alana harus menghadapi kesedihan seorang diri.

Beberapa hari setelah pemakaman, seorang pria kaya datang ke desa menggunakan mobil mewah. Pria tersebut bernama Arman. Ia adalah ayah kandung Alana yang selama ini tidak pernah hidup bersama mereka.

Arman kemudian membawa Alana pergi dari desa menuju kota.

Perjalanan tersebut menjadi awal kehidupan baru bagi Alana.

Untuk pertama kalinya, Alana melihat kehidupan yang sangat berbeda dari kehidupannya di desa. Rumah ayahnya sangat besar dan mewah. Halamannya luas, terdapat banyak kendaraan, serta para pekerja yang mengurus berbagai keperluan rumah.

Alana seharusnya merasa bahagia karena akhirnya dapat tinggal bersama ayah kandungnya.

Namun, kenyataan justru membuatnya semakin terluka.

Di dalam rumah tersebut, Alana tidak pernah benar-benar diperlakukan sebagai anak. Di hadapan banyak orang, Arman tidak pernah memperkenalkannya sebagai putrinya. Alana hanya dianggap sebagai anak dari seorang kenalan lama.

Perlakuan tersebut membuat Alana merasa asing di rumah ayahnya sendiri.

Ia memiliki kamar yang nyaman, makanan yang cukup, pakaian yang bagus, dan segala kebutuhan yang terpenuhi. Tetapi semua kemewahan itu tidak mampu menggantikan satu hal yang sangat ia inginkan, yaitu pengakuan dari ayahnya.

Kesedihan Alana semakin bertambah ketika ia harus pindah sekolah.

Sekolah barunya merupakan salah satu sekolah ternama di kota. Sebagian besar siswa berasal dari keluarga berada. Mereka terbiasa menggunakan barang-barang mahal dan menjalani kehidupan yang jauh berbeda dari kehidupan Alana sebelumnya.

Alana datang dengan penampilan sederhana dan kebiasaan yang masih sangat dipengaruhi kehidupan desa.

Perbedaan tersebut membuatnya menjadi perhatian.

Pada awalnya, Alana hanya mendapatkan tatapan aneh dari teman-temannya. Lama-kelamaan, tatapan tersebut berubah menjadi ejekan. Cara berbicaranya ditertawakan. Asal-usulnya dijadikan bahan olok-olok. Pakaian dan tas sederhananya dianggap memalukan.

Alana mulai sering dijauhi.

Tempat duduknya pernah dipindahkan. Buku-bukunya pernah disembunyikan. Beberapa barang miliknya juga pernah dirusak.

Setiap kejadian tersebut membuat hati Alana semakin terluka.

Namun ia memilih diam.

Alana tidak ingin membuat masalah. Ia juga tidak ingin membalas perlakuan buruk yang diterimanya.

Setiap kali pulang sekolah, ia hanya menyimpan semua kesedihan tersebut di dalam hati.

Di rumah pun ia tidak menemukan tempat untuk mencurahkan perasaannya.

Ayahnya selalu sibuk dengan pekerjaan dan kehidupan sosialnya. Hubungan mereka terasa seperti hubungan antara orang asing yang tinggal di bawah satu atap.

Alana akhirnya mulai terbiasa menyembunyikan air matanya.

Ia menangis diam-diam di dalam kamar ketika tidak ada seorang pun yang melihat.

Di tengah kesepian itu, Alana menemukan satu-satunya cara untuk mengubah hidupnya.

Ia belajar.

Setiap hari ia menghabiskan waktu dengan membaca buku dan mengerjakan tugas. Ia tidak ingin membuktikan sesuatu kepada orang-orang yang merendahkannya. Ia hanya ingin membuktikan kepada dirinya sendiri bahwa keadaan tidak menentukan masa depannya.

Prestasinya perlahan meningkat.

Guru-guru mulai memperhatikan kemampuan Alana. Ia sering mendapatkan nilai tinggi dan mulai dipercaya mengikuti berbagai kegiatan sekolah.

Anak yang sebelumnya dianggap sebagai anak desa mulai dikenal karena kecerdasannya.

Suatu hari, sekolah mengadakan perlombaan akademik. Alana menjadi salah satu peserta. Banyak orang tidak menyangka bahwa gadis sederhana yang sering dijauhi itu mampu bersaing dengan siswa-siswa terbaik.

Alana akhirnya berhasil meraih juara pertama.

Penghargaan tersebut menjadi bukti bahwa asal seseorang bukanlah ukuran kemampuan.

Kabar mengenai prestasi Alana akhirnya sampai kepada Arman.

Untuk pertama kalinya, Arman mulai menyadari betapa kuatnya putrinya.

Ia mulai memahami bahwa selama ini Alana tidak pernah meminta kekayaan. Alana hanya membutuhkan kasih sayang dan pengakuan sebagai seorang anak.

Namun, perubahan Arman tidak terjadi dalam satu malam.

Ia masih membutuhkan keberanian untuk menghadapi keluarganya dan orang-orang di sekitarnya.

Sementara itu, Alana terus menjalani hidupnya.

Ia tidak lagi terlalu memikirkan perkataan orang lain.

Perlahan, ia mulai menemukan beberapa teman yang benar-benar tulus menerimanya. Ia juga mulai berani berdiri untuk dirinya sendiri ketika mendapatkan perlakuan yang tidak baik.

Alana akhirnya memahami bahwa diam bukan berarti lemah.

Terkadang, seseorang memilih diam karena tidak ingin menjadi seperti orang yang telah menyakitinya.

Waktu terus berjalan.

Hubungan Alana dan ayahnya perlahan membaik. Arman mulai menunjukkan perhatian yang selama ini tidak pernah diberikan. Ia mulai menyadari kesalahannya dan berusaha memperbaiki hubungan dengan putrinya.

Pada sebuah acara sekolah yang dihadiri banyak orang, Arman akhirnya mengambil keputusan yang selama ini selalu ia hindari.

Di hadapan keluarga, guru, dan orang-orang yang mengenalnya, Arman mengakui Alana sebagai anak kandungnya.

Bagi orang lain, pengakuan tersebut mungkin terlihat sederhana.

Namun bagi Alana, hal itu merupakan sesuatu yang telah ia tunggu sejak pertama kali meninggalkan desa.

Air mata yang selama ini disimpan akhirnya jatuh.

Bukan karena kesedihan, tetapi karena luka yang selama ini ia pendam perlahan mulai menemukan tempat untuk sembuh.

Meskipun begitu, Alana tidak melupakan masa lalunya.

Ia tetap mengingat rumah kecil di desa.

Ia tetap mengingat ibunya.

Ia tetap mengingat masa ketika dirinya hanya memiliki seorang ibu dan kehidupan sederhana.

Bertahun-tahun kemudian, Alana kembali ke desa.

Ia berjalan menuju makam ibunya dengan membawa bunga.

Di hadapan makam tersebut, ia mengenang semua perjalanan yang telah dilaluinya.

Kehilangan ibunya.

Kepindahan ke kota.

Kehidupan bersama ayah yang kaya tetapi tidak mengakuinya.

Perundungan di sekolah.

Kesepian.

Tangisan yang disembunyikan.

Hingga akhirnya ia mampu berdiri dengan kedua kakinya sendiri.

Alana menyadari bahwa semua luka tersebut telah membuatnya menjadi seseorang yang lebih kuat.

Ia tidak lagi merasa malu karena berasal dari desa.

Ia tidak lagi merasa rendah karena pernah menjadi korban perundungan.

Ia juga tidak lagi membutuhkan pengakuan orang lain untuk mengetahui bahwa dirinya berharga.

Alana menatap langit sore dengan senyum kecil.

Kehidupannya memang tidak sempurna.

Namun ia telah belajar bahwa kehilangan seseorang tidak berarti kehilangan seluruh harapan.

Ia juga belajar bahwa kekayaan tidak selalu berarti kebahagiaan dan rumah besar tidak selalu menjadi tempat seseorang merasa diterima.

Kadang-kadang, rumah yang paling sederhana justru menyimpan kasih sayang yang paling besar.

Dan bagi Alana, rumah itu adalah kenangan tentang ibunya.

Seorang anak desa mungkin bisa diremehkan karena tempat ia berasal, tetapi tidak seorang pun berhak menentukan seberapa tinggi ia boleh bermimpi.

Alana akhirnya melangkah meninggalkan makam ibunya.

Ia tidak lagi berjalan sebagai anak kecil yang kehilangan arah.

Ia berjalan sebagai seorang gadis yang telah menemukan kekuatan dalam dirinya sendiri.

\jedahias
